%% GENERATED by selfunc_v411.py -- do not hand-edit.
\begin{table}
\centering
\scriptsize
\setlength{\tabcolsep}{4pt}
\caption{Selection function of the searched sample against a volume-limited reference census. Both columns are the same Gaia~DR3 selection --- $\varpi>0$, $\sigma_\varpi/\varpi\le0.1$ and $1/\varpi\le40$\,pc, 17\,566 objects --- of which the 89 searched stars are rows, so the two are selected alike by construction; and both are classified by one rule, in which BP$-$RP colour sets the class and the absolute G magnitude separates the degenerates that colour alone would place among the A~stars. The rule reaches 17\,278 census objects and 87 of the 89 searched stars, the rest carrying no Gaia BP or RP photometry. ``Ratio'' is the searched column fraction over the census column fraction; $>1$ is over-representation. The strongest selection effect here is distance, not colour.}
\label{tab:selfunc}
\begin{tabular}{@{}lrrr@{}}
\toprule
Property & Census & Searched & Ratio \\
\midrule
\multicolumn{4}{@{}l}{Distance} \\
\quad $d\le5$\,pc & 48 (0.3\%) & 12 (13.5\%) & 49.3 \\
\quad 5--10\,pc & 240 (1.4\%) & 13 (14.6\%) & 10.7 \\
\quad 10--20\,pc & 2069 (11.8\%) & 18 (20.2\%) & 1.7 \\
\quad 20--30\,pc & 5391 (30.7\%) & 19 (21.3\%) & 0.7 \\
\quad 30--40\,pc & 9818 (55.9\%) & 27 (30.3\%) & 0.5 \\
\midrule
\multicolumn{4}{@{}l}{Colour class$^{a}$} \\
\quad A & 61 (0.4\%) & 5 (6.3\%) & 17.4 \\
\quad F & 473 (2.8\%) & 8 (10.1\%) & 3.6 \\
\quad G & 978 (5.8\%) & 15 (19.0\%) & 3.3 \\
\quad K & 2221 (13.2\%) & 10 (12.7\%) & 1.0 \\
\quad M & 13\,082 (77.8\%) & 41 (51.9\%) & 0.7 \\
\midrule
\multicolumn{4}{@{}l}{Other axes} \\
\quad Confirmed planet hosts$^{b}$ & 26 (15.5\%) & 23 (25.8\%) & 1.7 \\
\quad White dwarfs$^{a}$ & 463 (2.7\%) & 8 (9.0\%) & 3.4 \\
\quad Searched stars / systems & -- & 89 / 82 & -- \\
\bottomrule
\end{tabular}

\smallskip
{\footnotesize $^{a}$Colour bins, not MK types. The boundaries are the colours at which the census's own temperature--colour relation crosses 7500, 6000, 5300 and 3900\,K: BP$-$RP $<0.36$ for A, 0.36--0.72 for F, 0.72--0.95 for G, 0.95--1.77 for K and $>1.77$ for M. Over the 8594 census objects that also carry a Gaia temperature the two classifications agree for 98 per cent, and every disagreement is between neighbouring bins. A star is counted as degenerate when it lies more than 4\,mag below the main sequence measured from the census itself: the two sequences are separated by 5 to 8\,mag here, so moving the cut by a magnitude either way changes the count by under 3 per cent. The five classes are normalised over the 79 non-degenerate searched stars and the 16\,815 non-degenerate census objects.\\ $^{b}$Both columns are the same positional join, of the 168-entry ALMA-covered work list and of its searched subset, to the NASA Exoplanet Archive \texttt{pscomppars} table accessed 2026 October 4. The 23 searched hosts carry 50 confirmed planets, 21 of the 23 surviving the deuterium-burning mass limit; the Gaia census carries no planet information, so the census column in this row is the work list and not the 17\,566-object Gaia census.}
\end{table}
